\subsection{Access changes in the Claude batch}
All fifteen judgments that changed between access conditions in the Claude batch concern the software and documentation criteria. For TruthfulQA and DesignQA in both repetitions, A1 and A2 move from 5 in the documents condition to 10 in the repository condition. In the documents condition those responses state that the implementation cannot be inspected and apply anchor 5; in the repository condition they describe the entry point or compare the implementation with the documented rule and apply anchor 10. For HumanEval, where the paper itself prints the estimator code, A2 is 10 in every session, and A1 is 5 in the first documents repetition and 10 in the second, so the same README produced both anchors within one configuration. No A1 or A2 judgment reaches 15, and each repository response explains why execution for the current edition is not documented. In the Codex batch the documents sessions had used 10 for A1 from a README command alone in four of six cases. The Claude judgments follow the anchor distinction between a claimed and an inspectable implementation more often, but this is a statement about adherence to anchor wording, not about the correctness of either rating.

The remaining changes concern A3, A4 and A5 and are not uniform. For DesignQA A3, the repository response reads CSV headers and confirms metadata fields that the documents response could not, which is a change supported by information only available in the repository. For TruthfulQA A3 the direction differs between repetitions, which indicates a sufficiency threshold about edition metadata applied to text present in both conditions. No measurement criterion changed with access in this batch. Repository files address what can be inspected; they did not alter the responses' reading of construct, checking, reliability, human comparison or exposure evidence.

\subsection{Repository findings that the paper does not contain}
Several repository responses record observations from inspection rather than from the paper. For HumanEval, the response reports that the README describes the execution call as commented out while the shipped file has it active, and that a README file name does not exist in the snapshot. For TruthfulQA, it reports an undocumented fixed reference answer appended to the true answer list before similarity scoring, and the absence of code for the 2025 binary setting. For DesignQA, it reports that the similarity metric in the paper is commented out in the shipped code and that the overall score divides by a fixed six regardless of the subsets supplied. The preparing assistant read the cited lines offline and found them as described. These are findings about the frozen snapshots, recorded without execution. The responses rate A2 as 10 rather than 0 in these cases, arguing that the main documented rule is matched and that the discrepancies concern reproducibility or edge cases. Whether a commented out metric that appears in a results table is a rule violation under anchor 0 is a judgment the rubric wording leaves open.

\subsection{Repeatability and scope choices}
Fifty five of sixty repeated judgments matched. The five changes are on A1, A3 and A5, and three of them concern A3, where the responses weigh partition statements and edition metadata differently between repetitions. For HumanEval B4, both conditions and both repetitions use 5 because the paper compares expected student performance with model results in its discussion. The Codex batch had used not applicable in the documents condition. For DesignQA B4, all four Claude responses use not applicable because no human performance is measured. These scope choices are stable within this batch, but the rubric does not say whether an interpretive comparison in a discussion section brings B4 into scope, so stability is not evidence that the choice is right.

One evidence entry in a DesignQA repository response cites the file listing itself as the source of an absence claim about a license file. The citation checker counts it as a source that does not exist, because the listing is a tool result rather than a file. The response states the inference plainly, and the snapshot has no license file. This shows that the mechanical source check and the support for a judgment can diverge in both directions.

The observations above cover all sixty judgments in the six preselected responses through the accompanying audit. The analyst is the same model as the reviewer, and no independent human analysis was performed. No semantic accuracy rate is inferred from quotation matching, from agreement between the batches or from the selected examples.
